\scanpage{42}
\begin{center}6 - 25 -\hfill n\textsuperscript{o} 392\end{center}

c)\unsure{.}Soient $k$ un corps ordonné, $a$, $b$ des éléments $>0$ de $k$ et $B$ l'anneau $k[X,Y]/(aX^2+bY^2+1)$. Montrer que $B$ est un anneau factoriel, (utiliser b) et l'exerc. 17 a) du \S\ 1).

d) Montrer que l'anneau $C=Q[X,Y]/(X^2+2Y^2+1)$ est factoriel, mais que l'anneau $C\otimes_Q Q(i)$, (avec $i^2=-1$), n'est pas factoriel.

7) Soient $K$ un anneau factoriel noethérien, $F$ un élément extrémal de l'anneau de polynômes $K[X_1,\ldots,X_n]$ ; on suppose que lorsque l'on attribue à chaque $X_i$ un poids $q(i)>0$, ($1\leq i\leq n$), $F$ est isobare et de poids $q>0$. Soit $A$ l'anneau engendré, dans une clôture algébrique $\Omega$ de $K(X_1,\ldots,X_n)$ par $K$, les $X_i$, ($1\leq i\leq n$), et une racine $z$ du polynôme $Z^c-F$, où $c$ est un entier étranger à $q$. Montrer que $A$ est factoriel dans les deux cas suivants :

1) on a $c\equiv 1$ (mod. $q$) ;

2) tout $K$-module projectif de type fini est libre, (ce qui a lieu par exemple lorsque $K$ est un corps, ou un anneau principal, ou un anneau local).

(Dans le premier cas, considérer l'anneau de fractions $A[1/z]$ ; montrer qu'il est factoriel et appliquer la prop. 3 du n\textsuperscript{o} 4. Dans le second cas, considérer un entier $d$ tel que $cd\equiv 1$ (mod. $q$), et soit $z'\in\Omega$ tel que $z=z'^{d}$ ; l'anneau $B=A[z']$ est factoriel en vertu du premier cas, et est un $A$-module libre. Considérer $B$ comme un anneau gradué en prenant $z'$ de poids $q$, chacun des $X_i$ de poids $cdq(i)$, et se ramener à prouver que pour deux éléments homogènes $u$, $v$ de $A$, l'idéal $Au\cap Av$ est principal, en utilisant l'exerc. 16 du \S\ 1 ; considérer enfin l'idéal $Bu\cap Bv$ dans $B$ et utiliser le chap. I, \S\ 3, n\textsuperscript{o} 6, prop. 12). En particulier, si $K$ est un corps, $a$, $b$, $c$ trois entiers $>0$ étrangers deux à deux, l'anneau $A=K[X,Y,Z]/(Z^a-X^b-Y^c)$ est factoriel.

8) Soient $A$ un anneau intègre, $x$, $y$, $z$ trois éléments non nuls de $A$ où $x$ est extrémal et $Ax\cap Ay=Axy$. Soit $S$ l'ensemble multiplicatif des $x^n$, ($n\geq 0$), et soit $B=S^{-1}A$ ; on considère les anneaux de séries formelles $A[[T]]$ et $B[[T]]$.

a) Soient $i$, $j$, $k$ trois entiers $\geq 0$ tels que $ijk-ij-jk-ki\geq 0$ et $z^i\in Ax^j+Ay^k$. On considère dans $A[[T]]$ l'élément $v=xy-z^{i-1}T$. Montrer qu'il existe un entier $t>0$ et une série $v'=y^tx^{-1}+b_1x^{-2}T+\ldots\ \ldots+b_{n-1}x^{-n}T^{n-1}+\ldots$ dans $B[[T]]$ tels que $vv'\in A[[T]]$, (on déterminera par récurrence les $b_n$, en procédant dans $N$ par intervalles de longueur $ij$ ;

\margin{41}

